\chapter{Implementation}\label{ch:implementation}
This chapter presents the requirements, selects a mapping library, introduces WebGL, and showcases the final implementation.

\section{Requirements and Design Choices}

The AirScout website will include a dedicated page for the map, which is used to visualize sensor data from all sensor boxes. This leads to the following requirements:
\begin{itemize}
    \item It must be trivial to understand the map and identify possibly interpolated values across an area. This can be achieved with a legend and appropriate color mapping.
    \item It must be performant enough to remain stable on low-end devices when there is a large number of data points.
    \item Areas with a lower density of data points must remain visible\\ (see\ \autoref{subsec:flaws}).
\end{itemize}

\section{Why WebGL}

Because of the computational complexity and the stated requirements, performance is a major consideration when evaluating implementation options. The main performance challenge comes from the fact that every visible pixel needs to consider every data point, or at least every data point within range. For this reason, the map is going to be implemented using a \emph{shader} (see\ \autoref{sec:webgl}). Shaders are programs that run on the GPU and are highly parallelized. Specifically, the map will be implemented as a \emph{fragment shader}, which is a type of shader that can calculate the color of every pixel in parallel.

To execute shaders in a web application, WebGL is the standard option. WebGL is a graphics API for the web based on OpenGL ES (see\ \autoref{sec:webgl}). It provides a JavaScript interface for rendering 2D and 3D graphics using the GPU\@. WebGL is supported on all modern browsers and is widely used for web-based graphics applications, such as games, data visualizations, and interactive maps.

\section{Mapping Libraries}

This section examines different JavaScript mapping libraries that can be used as a foundation for the map. They are evaluated based on the ease of adding a custom fragment shader as an overlay, the required development effort, rendering performance, and stability.

Used terminology:
\begin{itemize}
    \item Raster layers/tiles: map layers that are rendered as images. They are easier to work with and have better compatibility across browsers, but are less performant and flexible than vector layers.
    \item Vector layers/tiles: map layers that use vector graphics to render map features, similar to SVG\@. They are more performant and flexible than raster layers, but require more development effort and have lower compatibility across browsers.\ \cite{libs:tiles}
    \item Canvas: A 2D rendering context that allows for drawing graphics using JavaScript and HTML\@. It is easier to work with than WebGL, but is less performant and flexible.\\\cite{libs:canvas}
    \item GeoJSON\@: A format for encoding geographic data structures using JSON\@. It is widely supported and easy to work with.\ \cite{other:geojson}
\end{itemize}

\subsection{Mapbox GL JS}

\begin{quote}
    \textit{Mapbox GL JS is a client-side JavaScript library for building web maps and web applications with Mapbox's modern mapping technology. You can use Mapbox GL JS to display Mapbox maps in a web browser or client, add user interactivity, and customize the map experience in your application}\ \cite{libs:mapbox}.
\end{quote}

Mapbox GL JS (from now on referred to as \emph{Mapbox}) is built on top of Mapbox GL, a graphics library that relies on WebGL to render interactive maps directly within modern web browsers. Using WebGL, which provides hardware-accelerated graphics through the GPU, it is able to efficiently display large vector datasets, smooth animations, and highly interactive maps. This enables high performance even when handling dynamic data and frequent user interactions.

Although Mapbox provides great functionality and excellent performance, it is distributed under a commercial pricing model. A limited free-use tier is available, which currently permits up to 50,000 map loads per month (approximately 1,700 per day). Beyond this threshold, usage fees increase progressively depending on the number of map loads and selected services. It is unlikely that our project will reach this threshold, but if a larger school participated, reaching 1,700 requests would be trivial.

Displaying custom WebGL content in a Mapbox application is relatively straightforward. Developers can add a custom layer to the map style, which provides access to the underlying WebGL rendering context, as well as the transformation matrix that transforms the map’s coordinate system to display coordinates. Within this custom layer, shaders, buffers, and other necessary WebGL resources are set up. The layer must also define a render method that is called on each animation frame by the engine. The rendering method issues the required WebGL draw calls, allowing custom graphics to follow the camera movements of the map and integrate with the overall rendering pipeline.

\subsection{Leaflet}

Leaflet\cite{libs:leaflet} is an open-source JavaScript library designed to build interactive maps in web applications. In contrast to WebGL-based engines, Leaflet mainly uses HTML, CSS, and Canvas. This approach makes it lightweight and widely compatible across many browsers, while also remaining easy to learn and integrate. Its core functionality focuses on essential mapping features, such as tiled maps, markers, popups, vector layers, and event handling.

Because Leaflet does not use hardware-accelerated graphics by default, its performance is worse than that of WebGL-based solutions. For most use cases, Leaflet performs reliably and efficiently. However, when rendering very large datasets or highly dynamic visualizations, additional plugins or different rendering methods may be necessary to maintain acceptable performance. Because our data will be rendered in a shader and requires significant computation, the performance impact of Leaflet itself should be negligible.

One of Leaflet’s main advantages is its open-source license and extensive plugin ecosystem. As of 25.2.2026, Leaflet has support for 569 plugins, covering clustering, heatmaps, geospatial analysis, and integration with various tile providers. As a result, Leaflet is often used for academic projects and small- to medium-scale applications, and it is the most widely used JavaScript mapping library overall~\cite{libs:leaflet_usage}.

Extending Leaflet with custom rendering logic is possible through the use of plugins. This thesis will use Leaflet.TileLayer.GL~\cite{leaflet:tilelayergl}. This allows for the creation of custom layers that receive the underlying map texture and execute a fragment shader on each tile.

\subsection{Tangram}

Tangram~\cite{tangram:home} is an open-source WebGL-based mapping library created by Mapzen. Similar to Mapbox, it renders interactive and highly performant vector maps using the GPU\@.

A special feature of Tangram is its YAML- or JSON-based configuration\@. All aspects of the map are specified in a ``scene file'', including data sources, layers, different styling rules, custom shaders, and other rendering parameters. This provides a simple and easy-to-learn interface that also allows advanced configuration with shaders.

Tangram uses vector tiles by default and is designed for data from OpenStreetMap or Nextzen's vector tile service. The community has maintained the project since Mapzen's shutdown in 2018, and it is still actively developed.

Adding custom shaders to a Tangram map is possible by adding them to the scene configuration and using them in a custom layer. However, the documentation on custom shaders is sparse, and the API for them is not as well-defined as in Mapbox\ \cite{tangram:docs}.

\subsection{OpenLayers}

OpenLayers\ \cite{openlayers:main} is an open-source JavaScript mapping library similar to Mapbox. It supports many geospatial data formats and services such as raster tiles, vector data, and standards published by the Open Geospatial Consortium (OGC), such as WMS and WFS\@. Because of this, OpenLayers is often used in professional contexts where support for many geospatial data sources and formats is required. It is also used to power GIS applications, such as Pozi~\cite{openlayers:pozi} and yeymaps~\cite{openlayers:yeymaps}.

OpenLayers supports both raster and vector tiles and uses both Canvas and WebGL rendering. It is focused on displaying GIS data and building lightweight applications, which run well even on low-end and mobile devices.

Adding custom shaders to OpenLayers is possible, but the documentation on this topic is very limited. Our use case of rendering an overlay on top of map tiles using a custom fragment shader is not directly supported, and would need to be implemented by creating a custom layer and manually handling WebGL\ \cite{openlayers:githubgl}\@.

\subsection{amCharts}
amCharts~\cite{amcharts:home} is a commercial JavaScript library for creating data visualizations, such as charts, graphs, and maps. Unlike the other frameworks and libraries mentioned previously, amCharts is primarily focused on data presentation and visualization rather than interactive mapping. It provides components for statistical diagrams and has a dedicated module for maps (amCharts Maps), which is designed for GeoJSON data and supports various map projections.

While previous versions used SVGs, amCharts Maps version 5 uses Canvas for rendering, which improves performance with large datasets and allows for more complex visualizations (like heatmaps and hillshade, which cannot be broken down into individual objects)\@. amCharts is distributed under a dual licensing model. A free version is available as long as the amCharts branding is included in the visualizations. For commercial use or to remove the branding, a commercial license must be purchased.

Charts and maps are defined using the configuration system based on JavaScript objects. It defines chart types, data series, axes, legends, and interactive behaviors with configuration objects or API methods. For maps, geographic data is usually provided in the GeoJSON format or with predefined map files\@. amCharts manages rendering internally and does not provide an API for interacting with internal rendering. While amCharts supports many configuration options, such as for styling and animation, it is primarily meant for data representation rather than low-level graphics programming or direct manipulation of rendering contexts.

\subsection{Selection of a Mapping Library}

In conclusion:
\begin{itemize}
    \item Mapbox provides great performance and functionality, but the commercial pricing model is a major drawback.
    \item Leaflet is widely used and has a very large plugin ecosystem. It also has good support for custom shaders through the Leaflet.TileLayer.GL plugin.
    \item Tangram has a unique scene configuration and uses WebGL-based rendering, but lacks documentation for custom shaders.
    \item OpenLayers supports many geospatial data formats, but also suffers from limited documentation and functionality for custom shaders.
    \item amCharts is focused more on data visualization than interactive mapping and does not provide an API for custom shaders, making it unsuitable for our use case.
\end{itemize}

Based on the above analysis, Leaflet was chosen as the mapping library for this project. It is open-source, widely used, and supports a large number of plugins for different use cases. The Leaflet.TileLayer.GL by Iván Sánchez Ortega provides a simple way to add custom shaders as an overlay on top of map tiles, which is exactly what is needed for this project. The performance of Leaflet should be sufficient for our use case, as the overlay will be the primary performance bottleneck.

\section{Introduction to WebGL}\label{sec:webgl}

\begin{quote}
    \textit{WebGL is a technology that enables drawing, displaying, and interacting with sophisticated
interactive three-dimensional computer graphics [$\ldotp\!\ldotp\!\ldotp$] from within web browsers.}\ \cite{webgl:book}
\end{quote}

WebGL is a browser API that provides a JavaScript interface for rendering graphics using the GPU\@. It is based on OpenGL ES, a subset of the OpenGL graphics API designed for embedded systems. WebGL is developed and maintained by the Khronos Group, an organization of many companies that operate in the field of graphics and computing. The Khronos Group also develops other graphics APIs, such as the previously mentioned OpenGL and OpenGL ES, as well as Vulkan, a graphics API similar to OpenGL, SPIR-V, a binary intermediate language for shaders, and many others.\ \cite{webgl:khronos}

Two versions of WebGL are currently used: WebGL 1.0, which is based on OpenGL ES 2.0 and WebGL 2.0, which is based on OpenGL ES 3.0. This thesis will exclusively use WebGL 2.0, as it provides features that are essential for the implementation, such as support for data textures with floating-point values.\ \cite{webgl:2fundamentals}

\subsection{Rendering Pipeline Overview}

The \emph{rendering pipeline} is a series of steps, also called stages, that are executed in order to display graphics on the screen. The input to the pipeline is a set of vertices, which are points in three-dimensional space that make up the rendered object. Vertices can also have associated attributes such as color, texture coordinates, and normal vectors. The result of the pipeline is an image that is usually displayed on the screen. The rendering pipeline can be divided into the following stages:

\begin{itemize}
    \item Vertex Processing: Here, the vertex shader is executed on each vertex. It transforms the coordinates into normalized device coordinates, which later get converted to positions on the screen. It also passes the vertex attributes to the next stages.
    \item Primitive Assembly and Rasterization: The vertices are assembled into primitives, such as points, lines, or triangles. The rasterizer then converts the primitives to fragments, which are potential pixels. Each fragment also has attributes, similar to vertices, which are interpolated from the vertex attributes.
    \item Fragment Processing: In this stage, the fragment shader is executed for each fragment from the previous stage. It calculates the color and other attributes, such as opacity, of each fragment.
    \item Alpha Blending and Output: The final color of each fragment is blended with other objects in the scene based on their distance from the camera and their opacity. The resulting image is then displayed on the screen.
\end{itemize}

\ \cite{opengl:book}

\subsection{Textures}

A texture is an image that can be applied to the surface of a 3D object to give it color and detail. In WebGL, textures are usually used to store images, but they can also be used to store arbitrary data, such as the data points for our map. This is done by using a special type of texture called a \emph{data texture}, which can store floating-point values instead of colors\ \cite{webgl:mdn}.

\section{Getting Shaders onto Maps}

\subsection{Leaflet.TileLayer.GL}

Leaflet.TileLayer.GL is a plugin for Leaflet that was created by Iván Sánchez Ortega in 2016. It allows for the creation of custom layers, which take one or more map layers as input to a custom fragment shader.\ \cite{leaflet:tilelayergl}

The plugin works by creating a single canvas, which is used to render the map tiles. This canvas and its WebGL context are reused for all tiles, and the shader is executed separately for each tile. The result is then drawn on top of the tile using the Canvas 2D context.\ \cite{mdn:canvas2d}

The vertex shader used by the plugin cannot be customized, but it only sets the vertex positions and passes the texture coordinates to the fragment shader. The plugin also injects a fixed header into the fragment shader, which includes the GLSL version, varyings, and uniforms. This means that the fragment shader only needs to define the \verb|main| function and can access the provided parameters. With this approach, all user-defined uniforms need to be injected into the shader code.

Leaflet.TileLayer.GL initially used WebGL 1.0, but it was ported to WebGL 2.0 for this thesis. This required major changes to the shader code, because WebGL 2.0 uses GLSL version 3.00 ES, which is not compatible with GLSL version 1.00 used in WebGL 1.0\ \cite{webgl:2fundamentals}.

Leaflet.TileLayer.GL supports custom uniforms, but only floats and vectors up to length 4, because they can fit into a single variable type (float, vec2, vec3, or vec4). This is achieved by injecting the custom uniforms into the shader code. For example, the bandwidth uniform is declared with the following code in JavaScript:
\begin{lstlisting}[language=JavaScript,caption=Injecting Custom Uniforms into the Shader]
uniforms: {
    uBandwidth: 1.0,
},
\end{lstlisting}
which injects the following into the shader:
\begin{lstlisting}[language=GLSL]
uniform float uBandwidth;
\end{lstlisting}

\subsection{Changes to Leaflet.TileLayer.GL}\label{subsec:tlgl_changes}

For the map to work, many data points need to be passed to the shader, which is not possible with just regular uniforms. To solve this, the library was extended with support for array uniforms and data textures.

\subsubsection{Array Uniforms}

Array uniforms are a simple extension of regular uniforms. They use arrays of the basic types and require a \verb|size| parameter to specify the type of each entry. For example, the color stops are passed as an array of \verb|vec4|s, which is declared in JavaScript as follows:
\begin{lstlisting}[language=JavaScript,caption=Passing Color Stops as Array Uniforms]
arrayUniforms: {
    uColorStops: {
        size: 4,
        value: [...],
    },
},
\end{lstlisting}
which injects the following into the shader:
\begin{lstlisting}[language=GLSL,caption=Injected Array Uniforms]
uniform vec4 uColorStops[3];
\end{lstlisting}
WebGL does not allow arrays of length 0, so if the array is empty, the following macro is injected instead:
\begin{lstlisting}[language=GLSL,caption=Empty Array Macro]
#define UCOLORSTOPS_EMPTY
\end{lstlisting}

This allows for the use of a single shader program that can handle both cases. Note that values must be passed as a single flattened array of numbers.

This approach, however, has two major disadvantages. First, the shader needs to be recompiled if the length of the array changes. Second, the maximum length of the array is limited by the maximum number of uniform vectors that can be passed to the shader, and most devices support only 1024.\ \cite{webglsurvey:uniform_vectors}\@.

\subsubsection{Data Textures}

To solve the problems of array uniforms, data textures can be used. Data textures are regular textures that store data instead of colors. They can be accessed in the shader using the \verb|texelFetch| function, which retrieves the value of a specific texel (pixel) from the texture.

For example, the data points for the shader are defined in JavaScript as follows:
\begin{lstlisting}[language=JavaScript,caption=Passing Data Points as a Data Texture]
dataTextures: {
    uTexPoints: {
        format: gl => ({
            internalFormat: gl.RGB32F,
            format: gl.RGB,
            type: gl.FLOAT,
        }),
        size: 3,
        value: this.mapDataPoints(dataPoints),
    },
},
\end{lstlisting}
which injects the following into the shader:
\begin{lstlisting}[language=GLSL]
uniform sampler2D uTexPoints;
uniform int uTexPoints_len;
\end{lstlisting}

Empty textures are supported, so there is no need for a special case in the shader. An initial implementation attempt used a texture with height 1 and the same width as the number of data points. This is problematic because textures are limited by their maximum size. Most devices support textures up to 16384$\times$16384 pixels, but 2048$\times$2048 should be used for compatibility.\ \cite{webglsurvey:texture_size}.

The new approach uses a texture with the maximum available width and wraps to the next row when the number of data points exceeds the width. This allows for a minimum of around 4 million data points, but slightly increases the complexity of the shader code.

\begin{lstlisting}[language=GLSL,caption=Accessing Data Points from Data Texture]
int width = textureSize(uTexPoints, 0).x;
for (int i = 0; i < uTexPoints_len; i++) {
    vec3 point = texelFetch(uTexPoints, ivec2(i % width, i / width), 0).rgb;
}
\end{lstlisting}

\section{Implementation}

This section examines the final implementation of the shader, explains how data is passed, and describes all important nuances and differences from the reference implementation.

\subsection{Shader}

\begin{lstlisting}[language=GLSL,caption=Final Shader Implementation]
float kernel(float distance, float bandwidth) {
    float u = distance / bandwidth;
    if (u > 1.0) {
        return 0.0;
    }
    return (1.0 - u * u);
}

float linstep(float a, float b, float t) {
    return clamp((t - a) / (b - a), 0.0, 1.0);
}

void main(void) {
    // Classic texel look-up
    vec4 texelColor = texture(uTexture0, vTextureCoords.st);
    float luminosity = 0.299 * texelColor.r + 0.587 * texelColor.g + 0.114 * texelColor.b;
    const float greyscaleMix = 0.6;

    vec3 finalColor = mix(vec3(luminosity), texelColor.rgb, greyscaleMix);

    #if defined(UCOLORSTOPS_EMPTY)
    fragColor = vec4(finalColor, texelColor.a);
    #else
    finalColor = clamp(finalColor, 0.0, 1.0);

    const float PI = 3.14159265358979323846;
    float zoomScale = pow(2.0, uTileCoords.z);
    vec2 coords = (uTileCoords.xy + vTextureCoords.st);
    vec2 transformed = (coords * 2.0 * PI / zoomScale - PI) * vec2(1.0, -1.0);

    float lon = degrees(transformed.x);
    float lat = degrees(2.0 * atan(exp(transformed.y)) - (PI / 2.0));
    vec2 latLngCoords = vec2(lat, lon);

    float density = 0.0;
    float value = 0.0;
    float bandwidth = uBandwidth / pow(1.5, uTileCoords.z);
    int width = textureSize(uTexPoints, 0).x;

    for (int i = 0; i < uTexPoints_len; i++) {
        vec4 point = texelFetch(uTexPoints, ivec2(i % width, i / width), 0);
        vec2 dist = latLngCoords - point.yx;
        float distance = length(vec2(dist.x / cos(radians(latLngCoords.x)), dist.y));

        float dense = kernel(distance, bandwidth);
        density += dense;
        value += dense * point.z;
    }

    float pixelDensity = density / 10.0;
    pixelDensity = smoothstep(0.0, 1.0, pixelDensity * sqrt(uTileCoords.z / 5.0));
    float pixelValue = density > 0.0 ? value / density : 0.0;

    vec3 pixelColor = uColorStops[0].rgb;
    for (int i = 0; i < uColorStops.length() - 1; i++) {
        pixelColor = mix(
                pixelColor,
                uColorStops[i + 1].rgb,
                linstep(uColorStops[i].a, uColorStops[i + 1].a, pixelValue)
            );
    }
    fragColor = vec4(mix(finalColor, pixelColor, min(pixelDensity * 0.8, 0.7)), texelColor.a);
    #endif
}
\end{lstlisting}
\begin{itemize}
    \item \cdlines{1--7} define the kernel function. The shader uses the parabolic kernel because it cuts off at a distance of $u > 1$, allowing for further optimization without changing parameters.
    \item \cdlines{9--11} define the \verb|linstep| function that interpolates between two edges, similar to OpenGL's smoothstep\ \cite{opengl:esref}, but using linear instead of Hermite interpolation.
    \item \cdline{15} retrieves the original color of the map tile.
    \item \cdlines{16--19} apply a 40\% greyscale filter to improve the visibility of the overlay (60\% of the original color is retained).
    \item \cdlines{26--33} convert the pixel coordinates from the Web Mercator projection to latitude and longitude. The formula was derived from the Leaflet source code\ \cite{leaflet:source}\@.
    \item \cdline{37} adjusts the bandwidth based on the zoom level.
    \item \cdlines{40--48} loop over all data points and calculate the density and interpolated value for the current pixel using the kernel.
    \item \cdlines{50--52} calculate the local density, apply the \verb|smoothstep| function to it, and compute the interpolated value.
    \item \cdlines{54--61} calculate the color of the overlay using the provided color stops.
    \item \cdline{68} mixes the original tile color with the overlay color based on the density and outputs the final color.
\end{itemize}

\subsection{Implementation Details}\label{subsec:shader_details}

In detail, the shader does the following steps for each pixel:

First, the current bandwidth $h$ is calculated from the base bandwidth $B$ and the zoom level $z$.
\[
h = \frac{B}{1.5^{z}}
\]

Note that even though each zoom level doubles the resolution, the bandwidth is only reduced by a factor of $1.5$. This results in a smaller bandwidth at higher zoom levels, making finer details and individual data points visible at high zoom levels while keeping the overlay visible at low zoom levels.

Then, the pixel's coordinate $(x, y)$ is converted into its latitude $\varphi_{P}$ and longitude $\lambda_{P}$ using the inverse Web Mercator projection\ \cite{explorers:webmerc}:
\[
\begin{aligned}
\varphi_P &= 2 \cdot \arctan\left(\exp\left(\pi - \frac{2\pi y}{2^z}\right)\right) - \frac{\pi}{2}\\
\lambda_P &= \frac{2\pi x}{2^z} - \pi 
\end{aligned}
\]
The $x$ and $y$ coordinates are in global tile space, meaning that their integer part represents the coordinates of the tile and the fractional part represents the position within the tile.

The kernel for the NW estimator is defined as follows:
\[
\qquad K(\delta) = \max\left( 1-{\left( \frac{\delta}{h} \right)}^{2}, 0\right)
\]
where $\delta$ is the distance between the pixel and the data point.

The shader uses the parabolic kernel because of its sharp cutoff at $u > 1$. This allows for further optimizations by skipping data points that are too far away, and it provides a better visual appearance than the Gaussian kernel, which results in a broad, blurry overlay.

$d(i)$ is a shorthand for the weighting factor of the $i$-th data point. It is not explicitly defined in the shader, but is used here to simplify the mathematical notation. It is calculated using the kernel function and the distance between the pixel and the data point.
\[
d(i) = K\left(\sqrt{ {\left( \frac{\varphi_{P}-\varphi_{i}}{\cos(\varphi_{P})} \right)}^{2} + {(\lambda_{P}-\lambda_{i})}^{2}}\,\right)
\]

The distance is calculated using Euclidean distance in latitude and longitude, while accounting for the curvature of the Earth by dividing the latitude difference by the cosine of the pixel's latitude.

The shader then iterates over all $n$ data points to calculate the local density $\hat{\rho}$ and interpolated value $\hat{v}$.
\[
\begin{aligned}
\hat{\rho} &= \min\left(\frac{1}{10}\cdot\sqrt{ \frac{z}{5} \,}\sum_{i=1}^n d(i),1\right)\\
\hat{v} &= \frac{\sum_{i=1}^n d(i)\cdot v_{i}}{\sum_{i=1}^n d(i)}
\end{aligned}
\]

Note that the density is scaled by a factor of $\frac{1}{10}\cdot\sqrt{ \frac{z}{5} \,}$. This is done to increase the density at higher zoom levels, where points are sparser, and decrease it at lower zoom levels, where points are denser. The values $10$ and $5$ were found by trial and error and produce satisfactory visual results at all zoom levels using our test data. If large amounts of data points are inserted in the future, these constants may need to be adjusted, but they are sufficient for the current implementation. The density is also clamped to a maximum of $1$ to prevent overflow in the final opacity calculation. Unlike in the reference implementation, density is not divided by the total number of data points, which avoids the flaw that areas with a low density of data points can become invisible, as described in\ \autoref{subsec:flaws}.

The final opacity $\rho_{f}$ is calculated by applying the smoothstep function\ \cite{quilez:smoothstep} to further smooth the transition between low and high density. The output is then scaled by $0.8$ and clamped to a maximum of $0.7$ to prevent the overlay from becoming too opaque and obscuring the map tiles. This also renders high-density areas with a uniform opacity.
\[
\rho_{f} = \min\left(\max\left(0.8\cdot\left(3\hat{\rho}^2-2\hat{\rho}^3\right),0\right),0.7\right)
\]

Finally, the interpolated value is mapped to a color $c_f$ using linear interpolation between the color stops, and each RGB component is clamped to a maximum of 255. The color stops are defined as a list of tuples $(t_{i}, c_{i})$, where $t_{i}$ is the value of the color stop, $c_{i}$ is the corresponding color, and $n_c$ is the number of color stops.
\[
\begin{aligned}
i &= \max\left\{i \,|\,t_{i}\leq \hat{v}\right\}\quad i \in \{0,\,\dots\, ,n_{c}-1\}\\
c(t) &= \begin{cases}
c_{0} & \text{if} \, t\leq t_{0} \\
c_{n_c-1} & \text{if}\, t\geq t_{n_c-1} \\
c_{i}+\frac{t-t_{i}}{t_{i+1}-t_{i}}(c_{i+1}-c_{i}) & \mathrm{otherwise}
\end{cases}\\
c_{f} &= \mathrm{clamp}(c(\hat{v}), 0, 255)
\end{aligned}
\]
$i$ is the index of the color stop that is just below the interpolated value $\hat{v}$. $c(t)$ is the linear interpolation function, which is only used for the math representation, as the shader uses the built-in \verb|mix| function for linear interpolation. The final color is then calculated by mixing the original tile color with the overlay color based on the final opacity $\rho_{f}$\@.

\subsection{Passing Data}

The shader takes three inputs besides the underlying map tile: the base bandwidth $B$, the color stops, and the data points. Each one uses a different technique described in\ \autoref{subsec:tlgl_changes}.

\begin{itemize}
    \item The bandwidth is passed as a regular uniform, because it is a single float value that does not change often.
    \item The color stops are passed as an array uniform, because they are a small list of vec4 values that do not change often. The alpha component of the color stops is used to store the value of the color stop.
    \item The data points are passed as a data texture, because they are a large list of vec3 values that can change frequently. The red and green components store the latitude and longitude, and the blue component stores the value of the data point.
\end{itemize}

\begin{lstlisting}[language=JavaScript,caption=Passing Data to the Shader]
uniforms: {
    uBandwidth: 1.0,
},
arrayUniforms: {
    uColorStops: { size: 4, value: mappedColorStops },
},
dataTextures: {
    uTexPoints: {
        format: gl => ({
            internalFormat: gl.RGB32F,
            format: gl.RGB,
            type: gl.FLOAT,
        }),
        size: 3,
        value: mappedDataPoints,
    },
},
\end{lstlisting}

\verb|mappedColorStops| and \verb|mappedDataPoints| are both flattened arrays of numbers.